\section{Experimental Setup}
\label{sec:experimental-setup}

\subsection{Domains}
\label{sec:domains}

We evaluate on five tool domains spanning infrastructure operations and e-commerce, chosen to cover a range of graph topologies (Table~\ref{tab:domain-stats}).

\begin{table}[t]
\centering
\small
\begin{tabular}{l rrr r}
\toprule
\textbf{Domain} & \textbf{Tools} & \textbf{Types} & \textbf{T/T\%} & \textbf{Pruning} \\
\midrule
Kubernetes & 135 & 41 & 30\% & 63\% \\
Ansible & 108 & 40 & 37\% & 67\% \\
GitHub & 133 & 44 & 33\% & 56\% \\
CI/CD & 54 & 51 & 94\% & 87\% \\
Shopify & 170 & 61 & 35\% & 78\% \\
\bottomrule
\end{tabular}
\caption{Domain statistics. T/T\% = types-to-tools ratio. Pruning = average entity pruning (graph reachability constraint). CI/CD has near-parity types/tools yet achieves the highest pruning (87\%).}
\label{tab:domain-stats}
\end{table}

The domains vary along three dimensions relevant to our hypotheses:
\textbf{catalog size} (54--170 tools), \textbf{type density} (types-to-tools ratio from 30\% to 94\%), and \textbf{graph sparsity} (average entity pruning from 56\% to 87\%).
CI/CD is the most informative test case: with 51 types for 54 tools (94\% ratio), the type space provides almost no label reduction, yet achieves the highest entity pruning (87\%).
This shows that the graph constraint operates through topology, not label count.

Tool registries for Kubernetes, Ansible, GitHub, and CI/CD are constructed from their respective API documentation and operational patterns.
Shopify tools are derived from the Shopify Admin REST API, providing a non-infrastructure domain to test generalization.

\subsection{Models}
\label{sec:models}

We evaluate four language models spanning different scales and architectures:

\begin{itemize}
    \item \textbf{Granite 4.1 8B}: Open-weight, 8 billion parameters.
    \item \textbf{Qwen3-14B}: Open-weight, 14 billion parameters.
    \item \textbf{GPT-OSS-20B}: Open-weight, 20 billion parameters.
    \item \textbf{Claude Haiku 4.5}: Proprietary, parameter count undisclosed.
\end{itemize}

All models are used off-the-shelf without fine-tuning.
Our method should work with any model that can classify entity types from natural language, including models not trained for tool use.

\subsection{Query Categories}
\label{sec:query-categories}

Each domain includes 25--30 queries organized into six categories:

\begin{itemize}
    \item \textbf{Clean}: Well-structured queries with clear entity references (e.g., ``Get the logs for pods in the production namespace'').
    \item \textbf{Multi-hop}: Queries requiring 3+ tools in sequence (e.g., ``Show the resource usage of nodes running the cache deployment pods'').
    \item \textbf{Noisy}: Real-world language with irrelevant details (e.g., ``Something seems wrong with the API. Can you figure out where it is running?'').
    \item \textbf{Synonym}: Alternative terminology for entity types (e.g., ``workload'' for deployment, ``machine'' for host).
    \item \textbf{Ambiguous}: Queries where the target entity type is genuinely uncertain (e.g., ``Show me information about the frontend application'').
    \item \textbf{Multi-path}: Queries where multiple valid tool paths exist.
\end{itemize}

Total: 140 queries across 5 domains. Each query is annotated with ground-truth source and target entity types and the expected tool sequence (aligned to BFS shortest paths in the graph).

\subsection{Strategies}
\label{sec:strategies}

We compare six routing strategies:

\begin{enumerate}
    \item \textbf{Baseline}: The LLM selects tools directly from the full catalog presented in the prompt.
    \item \textbf{Retrieval}: Embedding-based retrieval returns the top-$k$ most similar tools; the LLM selects from this reduced set.
    \item \textbf{Graph}: Our method (typed composition search). The LLM predicts source and target entity types; BFS recovers the tool path.
    \item \textbf{Graph-Reverse}: Target-first variant with probability-based source selection using multiple completions.
    \item \textbf{Oracle}: Graph search with ground-truth entity types (upper bound on graph quality).
    \item \textbf{Model-Types}: Type prediction only, without tool routing (isolates type prediction accuracy).
\end{enumerate}

\subsection{Metrics}
\label{sec:metrics}

We report \textbf{precision}, \textbf{recall}, and \textbf{F1} computed over tool sets (predicted vs.\ expected).
We also report \textbf{hallucinated tool count}: tools returned by the model that do not exist in the registry.
For type prediction, we report source accuracy, target accuracy, and exact match (both correct).
